
% Language
%->8-------------------------------------------------------------------------------------------8<-%
\usepackage[brazilian]{babel} % Language is Brazilian Portuguese
\usepackage[utf8]{inputenc}   % File encoding is UTF-8
\usepackage[T1]{fontenc}      % Font encoding is T1
%->8-------------------------------------------------------------------------------------------8<-%

% Document Structure
%->8-------------------------------------------------------------------------------------------8<-%
\usepackage[a4paper, left=1.5cm, right=1.5cm, top=1.5cm, bottom=1.5cm]{geometry}
\usepackage{indentfirst}
\usepackage{changepage}
\usepackage{lmodern}
\usepackage[patch={none}]{microtype}
\newcommand{\frontmatter}{\cleardoublepage\pagenumbering{roman}} % Roman numeration in front-matter
\newcommand{\mainmatter}{\cleardoublepage\pagenumbering{arabic}} % Arabic numeration in main-matter
\newcommand{\backmatter}{\cleardoublepage\pagenumbering{arabic}} % Arabic numeration in back-matter
\setlength{\parskip}{\baselineskip}                              % Space between paragraphs
\newcommand{\nonport}[1]{\textit{#1}}                            % For texts in a foreign language
\newcommand{\algoname}[1]{\emph{#1}}                             % For algorithm names
\renewcommand{\labelenumi}{(\theenumi)}                          % Makes itemize be (number)
\newcommand{\fullref}[1]{\ref{#1} (\emph{\nameref{#1}})}
\usepackage{titlesec}
\titleformat{\section}
  {\normalfont\large\bfseries}
  {\thesection}{1em}{}
\titleformat{\subsection}
  {\normalfont\normalsize\bfseries}
  {\thesubsection}{1em}{}
\titleformat{\subsubsection}
  {\normalfont\normalsize\itshape}
  {\thesubsubsection}{1em}{}
%->8-------------------------------------------------------------------------------------------8<-%

% Math
%->8-------------------------------------------------------------------------------------------8<-%
\usepackage{amsmath, amssymb, mathdots, scalerel, mathtools, calcs}
\usepackage{amsthm}
\usepackage{thmtools}
\usepackage{titletoc}
\newcommand{\norm}[1]{\left\lVert#1\right\rVert}
\newcommand{\dotprod}[2]{\left\langle #1 \, \middle| \, #2 \right\rangle}
\renewcommand{\t}{^{\scriptscriptstyle\mathsf{T}\mkern-1mu}}
\newcommand{\st}{\,\,|\,\,}
\newcommand{\calciff}{\Updownarrow}

\declaretheoremstyle[
  headformat=\NAME~\NUMBER\let\thmcaption\NOTE\ifx\thmcaption\empty\else \textbf{\thmcaption}. \newline \fi,
  headpunct={},
  postheadspace=0.5em,
  headfont=\bfseries,
  notefont=\bfseries,
  bodyfont=\normalfont,
  spaceabove=0.3cm,
  spacebelow=0.3cm
]{definition-style}

\declaretheorem[
  style=definition-style,
  name=Definição,
]{definition}

\renewenvironment{proof}{
  \setlength{\parindent}{0pt}
  \setlength{\parskip}{0.6em}
  \noindent\textbf{Demonstração. } \par\nopagebreak
}{
  \qed
}

\declaretheoremstyle[
  headformat=\NAME~\NUMBER\let\thmcaption\NOTE\ifx\thmcaption\empty\else \textbf{\thmcaption}. \newline \fi,
  notefont=\bfseries,
  headpunct={},
  postheadspace=0.5em,
  headfont=\bfseries,
  bodyfont=\normalfont,
  spaceabove=0.3cm,
  spacebelow=0.3cm
]{theorem-no-title}

\declaretheorem[
  style=theorem-no-title,
  name=Teorema
]{theorem}

\declaretheorem[
  style=theorem-no-title,
  name=Lema
]{lemma}
%->8-------------------------------------------------------------------------------------------8<-%

% Code
%->8-------------------------------------------------------------------------------------------8<-%
  \usepackage{listings}
  \lstnewenvironment{code}[1][]{
    \lstset{
      basicstyle=\ttfamily,
      columns=fixed, % flexible
      breaklines=true,
      breakatwhitespace=true,
      frame=none,
      basewidth=0.5em,
      aboveskip=4pt,
      belowskip=4pt,
      mathescape=true,
      #1
    }
  }{}

  \newcommand{\var}[1]{\texttt{\detokenize{#1}}}
%->8-------------------------------------------------------------------------------------------8<-%

% Title
%->8-------------------------------------------------------------------------------------------8<-%
\title{Raciocício Equacional para Provar Equivalência entre Programas}
\author{João Victor Lopez Pereira}
\date{\today}
%->8-------------------------------------------------------------------------------------------8<-%

% Bibliography
%->8-------------------------------------------------------------------------------------------8<-%
\usepackage[backend=biber, style=abnt]{biblatex}
\usepackage{csquotes}
\addbibresource{repertoire.bib}
%->8-------------------------------------------------------------------------------------------8<-%
